\documentclass[12pt,a4paper]{report}

\usepackage[T1]{fontenc}
\usepackage[utf8]{inputenc}
\usepackage{lmodern}
\usepackage[top=1in,bottom=1in,left=1.25in,right=1in]{geometry}
\usepackage{setspace}
\usepackage{graphicx}
\graphicspath{{figures/}}
\usepackage{booktabs}
\usepackage{array}
\usepackage{tabularx}
\usepackage{longtable}
\usepackage{multirow}
\usepackage{amsmath}
\usepackage{amssymb}
\usepackage{xcolor}
\usepackage{listings}
\lstset{
  basicstyle=\ttfamily\small,
  breaklines=true,
  frame=single,
  numbers=left,
  numberstyle=\tiny\color{gray},
  keywordstyle=\bfseries\color{blue!70!black},
  commentstyle=\itshape\color{green!50!black},
  stringstyle=\color{red!60!black},
  showstringspaces=false,
  captionpos=b,
  language=Python
}
\usepackage[hidelinks]{hyperref}
\usepackage{fancyhdr}
\usepackage{titlesec}
\usepackage{caption}
\usepackage{subcaption}
\usepackage{float}
\usepackage{tikz}
\usepackage{pgfplots}
\pgfplotsset{compat=1.18}
\usetikzlibrary{shapes.geometric,arrows.meta,positioning,fit,backgrounds,calc}
\usepackage[style=ieee,backend=biber]{biblatex}
\addbibresource{references.bib}

\pagestyle{fancy}
\fancyhf{}
\fancyhead[L]{\small LeBlanc: Automated Prompt Enrichment and Iterative Repair}
\fancyhead[R]{\small \thepage}
\renewcommand{\headrulewidth}{0.4pt}

\titleformat{\chapter}[block]
  {\large\bfseries\centering}{Chapter \thechapter:}{0.5em}{}
\titlespacing*{\chapter}{0pt}{20pt}{10pt}

\onehalfspacing

% =========================================================
\begin{document}
\setcounter{page}{31}
\setcounter{chapter}{5}   % next \chapter is Chapter 6

% =========================================================
% CHAPTER 6 -- DISCUSSION
% =========================================================
\chapter{Discussion}
\label{ch:discussion}

\section{What the Results Show}
\label{sec:discmain}

The results across five research questions tell a consistent story: the LeBlanc
pipeline reduces vulnerability rates across all models, but the degree of
improvement depends heavily on the model's capacity to follow structured
instructions.

In plain mode, all three models produce fewer vulnerable runs than the Copilot
2022 baseline of 63.4\%. This is expected: Copilot was evaluated in 2022 against
a benchmark it was not trained to avoid, while our 2025 models have been
instruction-tuned with substantially more safety-relevant feedback. The
plain-mode results confirm that the field has moved, but they do not tell us
whether that movement is enough for production use. A 45--57\% vulnerability rate
in plain mode still means that roughly one in two generated code samples contains
at least one medium-severity finding.

Enrichment improves the picture for Command-A substantially (PEG = +42.9\%) and
has no measurable effect on Gemini 2.5 Flash (PEG = 0\%). The repair loop
is where the largest gains occur: all three models drop by at least 12 percentage
points in VR from enriched mode to repair mode, and Command-A reaches a
vulnerability rate of 18.2\% in full-pipeline mode. The practical implication
is that the repair loop contributes more than enrichment for Gemini 2.5 Flash,
and roughly equally for Command-A.

\section{The Llama 3.1-8B Enrichment Degradation}
\label{sec:llamadeg}

The most striking individual finding is Llama 3.1-8B's negative PEG of $-$33.3\%:
adding CWE warnings to the prompt causes the model to produce \emph{more}
vulnerable code than it does without them. This result requires careful
interpretation before drawing any conclusion about prompt enrichment as a strategy.

The 8B parameter size matters here. The enrichment block adds 80--120 tokens of
structured security requirements to the prompt. For a frontier model (Command-A,
Gemini 2.5 Flash) with strong instruction-following capability, these tokens
provide useful signal. For a model at the 8B scale, the additional tokens appear
to compete with or overshadow the core coding instruction. The model may be
attempting to address the security requirements and in doing so producing more
complex---and more vulnerable---code than it would with a simpler prompt.

This interpretation is consistent with published observations about small-model
degradation under complex prompts. It is \emph{not} evidence that CWE-aware
enrichment is counterproductive. The correct conclusion is narrower: CWE-aware
enrichment as implemented in Engine~A is not effective for models below a certain
instruction-following threshold, and Llama 3.1-8B is below that threshold.

A reviewer who observes this result and concludes that the enrichment strategy
is flawed is making a category error. The appropriate response is to control for
model capability, which is exactly what RQ5 does. The plan to re-run on Llama
3.3-70B---a model with the same architecture family but at a scale known to
follow complex instructions reliably---will provide a direct test of whether the
degradation is a function of the enrichment strategy or of the model size.

\section{Why Deserialization Does Not Converge}
\label{sec:discdeserial}

The CR@3 of 0\% for Command-A and Llama 3.1-8B on deserialization prompts
(Table~\ref{tab:category}) is the clearest ceiling effect in the dataset.
CWE-502 (Deserialization of Untrusted Data) typically manifests in Python as a
\texttt{pickle.loads} call on user-supplied data. The secure fix is not a
code-level change to the existing call but a design-level change: replace pickle
with a data-format library that does not allow arbitrary code execution (e.g.\
\texttt{json}, \texttt{msgpack}), or add an allowlist-based deserialiser.

Engine~C's repair prompt lists the finding and asks the LLM to fix the
vulnerability. For deserialization, the LLM's typical response is to add a
comment warning about the risk, or to wrap the call in a try/except block that
catches exceptions---neither of which removes the CWE-502 finding, and both of
which are therefore correctly flagged by Bandit on the re-scan. The model is not
producing incorrect code; it is producing code that addresses the symptom rather
than the cause.

This is a genuine limitation of the repair loop's current prompt structure. A
prompt that provides the model with an explicit replacement pattern (``replace
\texttt{pickle.loads} with \texttt{json.loads} or a safe equivalent'') would
likely achieve higher convergence. The current prompt does not do this; it
reports the finding and instructs the model to fix it without specifying
\emph{how}. This is an identified improvement for the next iteration of Engine~C.

\section{Gemini 2.5 Flash: Enrichment Insensitivity}
\label{sec:discgemini}

Gemini 2.5 Flash's PEG of 0\% indicates that its output distribution is not
affected by the CWE warning injection in Engine~A. Two explanations are
plausible. First, Gemini 2.5 Flash may already incorporate security-relevant
considerations in its base instruction tuning to a degree that makes the
explicit CWE warnings redundant. If the model already avoids CWE-89 on SQL
prompts without being told to, adding a warning about CWE-89 provides no
additional signal.

Second, the model's output may be determined more by its temperature setting
and top-$p$ sampling than by the additional tokens in the enrichment block.
At the low temperature of 0.2 used across all models, Gemini 2.5 Flash's
outputs are already near-deterministic, and the small perturbation from the
enrichment block may not shift the probability mass enough to change the output.

Distinguishing between these explanations would require a targeted experiment:
run the enrichment ablation at higher temperatures (0.7--1.0) to increase output
variance and check whether PEG becomes non-zero. This is out of scope for the
current evaluation but is a clean experimental question for future work.

\section{Limitations}
\label{sec:limitations}

Four limitations are worth stating explicitly.

\textbf{Single repetition.} Each prompt was run once per model per mode. LLM
outputs are stochastic: the same prompt can produce both a vulnerable and a clean
result across repeated runs. Single-rep results are directionally valid but
cannot support confidence intervals or variance estimates. Three repetitions per
prompt per model per mode (totalling approximately 900 runs) would be the minimum
for publication-quality claims.

\textbf{Bandit-only scanning.} The current implementation uses Bandit as the
sole scanner. Semgrep was included in the original design but removed from the
production configuration during development. Bandit covers a wide range of
common Python weaknesses but has known gaps: CWE-22 (Path Traversal) is caught
inconsistently depending on how the file path is constructed, and certain
logic-level access control weaknesses (CWE-285, CWE-306) may not be caught at
all if they manifest as missing checks rather than insecure calls. A dual-scanner
configuration (Bandit + Semgrep with the \texttt{p/python-security} rule pack)
would provide more complete coverage and should be restored in the next iteration.

\textbf{Dataset size and CWE coverage.} Thirty-three prompts across ten CWE types
gives an average of 3.3 prompts per CWE type. Per-CWE analysis at this granularity
has high variance. The deserialization category in particular has only two prompts,
making the 0\% CR@3 figure suggestive but not definitive. Expanding to the full
LLMSecEval set (120 prompts) or integrating the CodeSecEval benchmark would
provide more statistically reliable per-category estimates.

\textbf{False positive rate.} Bandit produces false positives, particularly for
medium-severity findings that are context-dependent. We did not perform a manual
verification pass on the scan results in this evaluation. An estimated false
positive rate of 10--15\% is consistent with published Bandit evaluations. If
10\% of our VULN results are false positives, the true vulnerability rates in
Table~\ref{tab:vr} would be approximately 5--6 percentage points lower than
reported. This does not reverse any of the directional findings but would reduce
the absolute VR figures.

% =========================================================
% CHAPTER 7 -- CONCLUSION AND FUTURE WORK
% =========================================================
\chapter{Conclusion and Future Work}
\label{ch:conclusion}

\section{Summary}
\label{sec:summary}

This report presented LeBlanc, a three-engine pipeline for automated prompt
enrichment, static vulnerability scanning, and iterative repair of LLM-generated
Python code. The system was evaluated across three models (Cohere Command-A,
Google Gemini 2.5 Flash, Meta Llama 3.1-8B), three pipeline modes (plain,
enriched, enriched with repair), and 33 prompts from the LLMSecEval benchmark,
producing 321 total runs logged to a reproducible SQLite database.

The main findings are:

\begin{itemize}
  \item All three models produce fewer vulnerable runs than the Copilot 2022
        baseline (63.4\%) in plain mode, confirming that LLM code security has
        improved since 2022 but has not reached acceptable levels.

  \item CWE-aware prompt enrichment reduces Command-A's vulnerability rate by
        42.9\% relative (PEG = +42.9\%) and has no statistically meaningful
        effect on Gemini 2.5 Flash (PEG = 0\%). Llama 3.1-8B degrades under
        enrichment (PEG = $-$33.3\%), a result we attribute to insufficient
        instruction-following capacity at 8B parameters.

  \item The iterative repair loop resolves vulnerabilities in 77\% of the 61
        runs that enter Engine~C, with a mean of 1.4 iterations. Command-A is
        the most repair-responsive model (CR@3 = 89.5\%); Llama 3.1-8B is the
        least (CR@3 = 62.5\%).

  \item Deserialization (CWE-502) is the hardest category to repair: the loop
        does not converge within three iterations for Command-A or Llama 3.1-8B.
        Injection, Access Control, and File Handling converge reliably.

  \item In full-pipeline mode (enriched + repair), Command-A achieves a
        vulnerability rate of 18.2\%, representing a 71.3\% improvement over
        the Copilot baseline.
\end{itemize}

LeBlanc is the first system to combine automated CWE-aware prompt enrichment,
multi-model parallel evaluation, and iterative automated repair with re-scanning
in a single deployable pipeline. The combination of a working web application and
a logged reproducible dataset contributes both a research tool and empirical data
to the open literature.

\section{Future Work}
\label{sec:futurework}

\subsection{Larger Open-Weight Models}
\label{ssec:fw_llama}

The most immediate next step is to re-run the full evaluation on Llama 3.3-70B
using the same Groq API access. The model identifier change in
\texttt{llm\_client.py} is a single line:

\begin{lstlisting}[language={}, caption={Model string replacement for Llama 3.3-70B}]
"llama3.3-70b": {
    "provider": "groq",
    "api_model": "llama-3.3-70b-versatile"
}
\end{lstlisting}

Llama 3.3-70B uses the same architecture family as Llama 3.1-8B but at a scale
that has demonstrated strong instruction-following on complex prompts. If its PEG
is positive---as expected from the model-size analysis in
Section~\ref{sec:llamadeg}---this would confirm that the enrichment degradation
observed in this evaluation is a size effect rather than a strategy effect, and
would establish Llama 3.3-70B as a viable open-weight option for deployment.

\subsection{Three-Repetition Dataset}
\label{ssec:fw_reps}

All current results are based on a single run per prompt per model per mode.
A three-repetition run ($33 \times 3 \times 3 \times 3 = 891$ runs) would
allow variance estimation across the vulnerability rate metrics. The
\texttt{run\_batch.py} script already supports \texttt{--reps N} for this
purpose. The cost is approximately 900 API calls; at current Groq and Cohere
pricing this is within budget for a semester project.

\subsection{Dual Scanner: Bandit and Semgrep}
\label{ssec:fw_semgrep}

The production Engine~B currently uses Bandit only. Semgrep with the
\texttt{p/python-security} rule pack covers additional CWE types---notably
CWE-22 (Path Traversal) via taint analysis---and would reduce the false negative
rate on path traversal and injection categories. The \texttt{run\_semgrep}
function exists in the codebase and ran correctly during development; it was
disabled for performance reasons. With the \texttt{p/python-security} rule pack
replacing \texttt{--config=auto}, startup time is acceptable and re-enabling
Semgrep should be straightforward.

\subsection{Engine C Repair Prompt Specialisation}
\label{ssec:fw_enginec}

The current repair prompt is generic: it lists findings and instructs the model
to fix them. For deserialization vulnerabilities (CWE-502) this generic
instruction does not produce convergence because the correct fix requires a
design-level change that the model does not infer from the finding description
alone. A set of category-specific repair prompt templates---one per
vulnerability category---would provide the model with explicit fix patterns
(``replace \texttt{pickle.loads} with \texttt{json.loads}'', ``use
\texttt{os.path.realpath} and validate the result against a base directory'')
and is expected to substantially improve CR@3 on the categories that currently
show low convergence.

\subsection{VS Code Extension}
\label{ssec:fw_vscode}

The web application is a research and demonstration platform. The intended
production deployment is a VS Code extension that calls the same backend
pipeline on the currently open file, highlights vulnerable lines in-editor,
and surfaces repair suggestions inline. The backend API design---with
\texttt{POST /api/run} accepting arbitrary prompt text and returning structured
scan and repair results---was specifically designed to support this integration.

The extension would be implemented in TypeScript using the VS Code Extension
API. On file save or manual trigger, it would extract the current file content,
send it as a prompt to \texttt{POST /api/run} in \texttt{enriched\_repair} mode,
and use the returned \texttt{scan\_results} to create diagnostic markers on the
affected lines. The \texttt{final\_code} from the repair result would be offered
as a one-click inline fix.

This work is out of scope for Semester VI but is the planned Semester VII
continuation. The API contract is already stable; no backend changes are required
to support the extension.

\subsection{Expanded Dataset and CWE Coverage}
\label{ssec:fw_dataset}

The 33-prompt LLMSecEval subset covers ten CWE types. The full LLMSecEval
benchmark has 120 prompts and broader CWE coverage. Integration of the
CodeSecEval benchmark (44 CWE types, 180 samples) would allow per-CWE analysis
with statistical power sufficient for publication. Both datasets are publicly
available and compatible with the existing \texttt{prompts.json} format.

\subsection{False Positive Rate Validation}
\label{ssec:fw_fpr}

A manual review of a 10\% random sample of scan findings would produce a false
positive rate estimate for the paper. At 321 runs with an average of approximately
2 findings per vulnerable run, a 10\% sample is roughly 64 findings. Each finding
requires a human to read the generated code and determine whether the Bandit
finding is a genuine vulnerability or a false alarm. This is achievable in a
single working day and would substantially strengthen the methodology section of
any conference submission.

\section{Closing Remarks}
\label{sec:closing}

LeBlanc demonstrates that a three-engine pipeline---automated enrichment,
deterministic static analysis, iterative repair---can systematically reduce the
vulnerability rate of LLM-generated Python code below the Copilot 2022 baseline
for every model tested. The system is deployable, the dataset is reproducible,
and the results are expressed in metrics that map directly to the five research
questions stated at the outset.

The most important open question the results raise is not whether the pipeline
works---the 77\% repair success rate and the 71.3\% RI for Command-A indicate
that it does---but how far these gains transfer to a wider range of models,
more CWE types, and higher-stakes code generation scenarios. The next evaluation
sprint, with Llama 3.3-70B, three repetitions, dual scanners, and category-aware
repair prompts, will answer that question directly.

% =========================================================
% APPENDIX A -- CWE Category Mapping
% =========================================================
\appendix
\chapter{CWE Category Mapping}
\label{app:cwemap}

Table~\ref{tab:cwemap} lists all CWE identifiers used in the LeBlanc pipeline,
their names, and the category each is assigned in \texttt{cwe\_categories.py}.

\begin{longtable}{@{}lll@{}}
\caption{CWE to category mapping used throughout the system}
\label{tab:cwemap} \\
\toprule
\textbf{CWE ID} & \textbf{CWE Name} & \textbf{Category} \\
\midrule
\endfirsthead
\multicolumn{3}{c}{\small\itshape (Table~\ref{tab:cwemap} continued)} \\
\toprule
\textbf{CWE ID} & \textbf{CWE Name} & \textbf{Category} \\
\midrule
\endhead
\bottomrule
\endlastfoot
CWE-20  & Improper Input Validation              & Injection \\
CWE-78  & OS Command Injection                   & Injection \\
CWE-79  & Cross-Site Scripting (XSS)             & Injection \\
CWE-89  & SQL Injection                          & Injection \\
CWE-94  & Code Injection                         & Injection \\
CWE-256 & Plaintext Storage of a Password        & Credential Management \\
CWE-259 & Use of Hard-coded Password             & Credential Management \\
CWE-327 & Use of a Broken or Risky Crypto Algorithm & Credential Management \\
CWE-328 & Use of Weak Hash                       & Credential Management \\
CWE-338 & Use of Cryptographically Weak PRNG     & Credential Management \\
CWE-521 & Weak Password Requirements             & Credential Management \\
CWE-798 & Use of Hard-coded Credentials          & Credential Management \\
CWE-285 & Improper Authorisation                 & Access Control \\
CWE-287 & Improper Authentication                & Access Control \\
CWE-306 & Missing Authentication for Critical Function & Access Control \\
CWE-384 & Session Fixation                       & Access Control \\
CWE-614 & Sensitive Cookie Without Secure Flag   & Access Control \\
CWE-732 & Incorrect Permission Assignment        & Access Control \\
CWE-200 & Exposure of Sensitive Information      & Information Exposure \\
CWE-209 & Error Message Contains Sensitive Info  & Information Exposure \\
CWE-532 & Insertion of Sensitive Info into Log   & Information Exposure \\
CWE-377 & Insecure Temporary File                & File Handling \\
CWE-434 & Unrestricted Upload of Dangerous File  & File Handling \\
CWE-22  & Improper Limitation of Pathname (Path Traversal) & Path Traversal \\
CWE-23  & Relative Path Traversal                & Path Traversal \\
CWE-36  & Absolute Path Traversal                & Path Traversal \\
CWE-502 & Deserialization of Untrusted Data      & Deserialization \\
CWE-611 & XML External Entity (XXE) Injection    & Deserialization \\
\end{longtable}

% =========================================================
% APPENDIX B -- Bandit Rule ID to Category Mapping
% =========================================================
\chapter{Bandit Rule to Category Mapping}
\label{app:banditmap}

Table~\ref{tab:banditmap} lists the Bandit rule IDs used by Engine~B's category
mapping function and their assigned categories.

\begin{longtable}{@{}lll@{}}
\caption{Bandit rule IDs and category assignments}
\label{tab:banditmap} \\
\toprule
\textbf{Rule ID} & \textbf{Rule Name} & \textbf{Category} \\
\midrule
\endfirsthead
\multicolumn{3}{c}{\small\itshape (Table~\ref{tab:banditmap} continued)} \\
\toprule
\textbf{Rule ID} & \textbf{Rule Name} & \textbf{Category} \\
\midrule
\endhead
\bottomrule
\endlastfoot
B102 & exec\_used                                & Injection \\
B307 & eval\_used                                & Injection \\
B404 & import\_subprocess                        & Injection \\
B602 & subprocess\_popen\_with\_shell\_equals\_true & Injection \\
B603 & subprocess\_without\_shell\_equals\_true   & Injection \\
B605 & start\_process\_with\_a\_shell            & Injection \\
B608 & hardcoded\_sql\_expressions               & Injection \\
B105 & hardcoded\_password\_string               & Credential Management \\
B106 & hardcoded\_password\_funcarg              & Credential Management \\
B107 & hardcoded\_password\_default              & Credential Management \\
B303 & md5 (weak hash)                           & Credential Management \\
B311 & random (insecure PRNG)                    & Credential Management \\
B324 & hashlib (weak algorithm)                  & Credential Management \\
B501 & request\_with\_no\_cert\_validation       & Credential Management \\
B104 & hardcoded\_bind\_all\_interfaces          & Access Control \\
B110 & try\_except\_pass                         & Information Exposure \\
B112 & try\_except\_continue                     & Information Exposure \\
B103 & set\_bad\_file\_permissions               & File Handling \\
B108 & hardcoded\_tmp\_directory                 & File Handling \\
B301 & pickle (deserialization)                  & Deserialization \\
B302 & marshal (deserialization)                 & Deserialization \\
B506 & yaml\_load (unsafe deserialization)       & Deserialization \\
B313--B320 & xml\_bad\_* (XXE variants)          & Deserialization \\
\end{longtable}

\printbibliography[heading=bibintoc,title={References}]

\end{document}
